\chapter{Validation}
\label{ch:validation}

This chapter asks whether the implementation in
Chapter~\ref{ch:implementation} meets the functional, non-functional, and KPI
targets set out in Chapter~\ref{ch:requirements} (Table~\ref{tab:kpis}). The
evidence comes from automated API tests, RAG evaluation harness runs, and
manual browser scenarios. Section~\ref{sec:critical-discussion} discusses how
much confidence each class of result actually warrants.

\section{Validation Plan}

Validation was organised at four levels. Unit tests cover pure helper functions.
Integration tests exercise the API end-to-end against a live PostgreSQL and Redis
stack. Browser scenarios on the Next.js client catch flows that are awkward to
assert at the HTTP layer alone. A dedicated RAG evaluation harness then measures
retrieval and answer quality with automated judge metrics when an API key is
available.

Relative to the solutions surveyed in Chapter~\ref{ch:state-of-the-art}, the
checks target capabilities that generic document stores or chatbots usually omit:
\textbf{department-scoped hybrid retrieval}, \textbf{cited streaming answers},
\textbf{per-user feature restrictions}, and \textbf{administrative audit trails}.
Unit coverage is kept deliberately narrow. The main automated quality gate is the
\textbf{74-test integration suite}, with manual UX checks filling in security and
restriction flows.

\section{Protocols and Experimental Conditions}

Table~\ref{tab:validation-setup} summarises the experimental environment.
Integration tests run against a seeded database with Docker-hosted PostgreSQL and
Redis. RAG evaluation uses a quick mode (optimiser-only, 34 cases) and a full
judge pipeline when a Gemini API key is configured. Manual scenarios were run in
the browser on the local development server; they appear in
Table~\ref{tab:manual-ux-tests}, with screenshot evidence in
Section~\ref{sec:manual-ux-validation}.

\begin{table}[H]
\centering
\caption{Experimental setup for validation}
\label{tab:validation-setup}
\footnotesize
\renewcommand{\arraystretch}{1.35}
\setlength{\tabcolsep}{4pt}
\begin{tabularx}{\textwidth}{|>{\RaggedRight\arraybackslash}p{3.0cm}|Y|}
\hline
\textbf{Parameter} & \textbf{Value} \\
\hline
Workstation & Windows 10/11; Node.js 20+ \\
\hline
Test framework & Vitest 3.2 + Supertest (\texttt{npm run test:api}) \\
\hline
Database / queue & PostgreSQL 16 + pgvector; Redis 7 (Docker Compose) \\
\hline
Test corpus & Seeded users, departments, and documents (\texttt{npm run db:seed}) \\
\hline
RAG evaluation & \texttt{eval:rag:quick} (34 topic cases); full judge if API key configured \\
\hline
Automated run date & 4 July 2026 (74/74 passed, duration $\approx$39\,s) \\
\hline
Manual UX & Browser on \texttt{localhost:3000}; scenarios T1--T5 (Table~\ref{tab:manual-ux-tests}) \\
\hline
\end{tabularx}
\end{table}

\section{Presentation of Results}

\subsection{KPI Validation Results}

Table~\ref{tab:kpi-results} maps each KPI from Chapter~\ref{ch:requirements} to its
measured outcome and validation conclusion.

\begin{table}[H]
\centering
\caption{KPI validation results}
\label{tab:kpi-results}
\footnotesize
\renewcommand{\arraystretch}{1.35}
\setlength{\tabcolsep}{3pt}
\begin{tabularx}{\textwidth}{|>{\RaggedRight\arraybackslash}p{2.2cm}|Y|>{\RaggedRight\arraybackslash}p{1.5cm}|>{\RaggedRight\arraybackslash}p{2.0cm}|>{\Centering\arraybackslash}p{1.4cm}|}
\hline
\textbf{KPI} & \textbf{Measurement method} & \textbf{Target} & \textbf{Obtained} & \textbf{Status} \\
\hline
Ingest reliability &
Uploads reaching ready state without failure &
$\geq 95\%$ &
100\% (14 doc.\ tests) &
Passed \\
\hline
Integration pass rate &
Full API integration suite &
74 / 74 &
\textbf{74 / 74} &
Passed \\
\hline
RAG topic (no API key) &
Quick eval optimiser pass rate &
$\geq 15/34$ &
\textbf{15 / 34} &
Passed \\
\hline
RAG judge (with API key) &
Automated faithfulness / relevance judge &
Documented &
Pending full run &
Partial \\
\hline
Security: refresh replay &
Rotated token revokes family &
Auth suite &
Pass (11 tests) &
Passed \\
\hline
Security: avatar CORP &
Cross-origin avatar embedding permitted &
Avatars test &
Pass (1 test) &
Passed \\
\hline
\end{tabularx}
\end{table}

\textbf{Analysis.} Five of the six KPI rows passed fully. On the seeded corpus,
ingest reliability clears the 95\% bar. The integration suite gives solid evidence
that API contracts match the design sequences in
Chapter~\ref{ch:architecture}. RAG topic classification hits the degraded-mode
threshold (15/34) without an external API key; full judge scores stay
\textbf{partial} until a complete run with credentials is available.

\subsection{Integration Test Results by Sprint Domain}

On 4 July 2026 the automated API suite reported \textbf{74 tests passed}
across 10 files with no failures. Table~\ref{tab:test-summary} groups the results by
sprint domain, in line with Table~\ref{tab:requirements}.

\begin{table}[H]
\centering
\caption{Validation test summary by sprint domain}
\label{tab:test-summary}
\footnotesize
\renewcommand{\arraystretch}{1.35}
\setlength{\tabcolsep}{3pt}
\begin{tabularx}{\textwidth}{|>{\RaggedRight\arraybackslash}p{0.9cm}|Y|>{\RaggedRight\arraybackslash}p{1.8cm}|>{\RaggedRight\arraybackslash}p{1.6cm}|>{\Centering\arraybackslash}p{1.3cm}|}
\hline
\textbf{Spr.} & \textbf{Objective} & \textbf{Expected} & \textbf{Obtained} & \textbf{Status} \\
\hline
S1 & Secure login and token refresh & 200 + new token & Pass (11 tests) & Passed \\
\hline
S2 & Document access control & 403/404 when restricted & Pass (14 tests) & Passed \\
\hline
S3 & Search and Ask restrictions & AI flag enforced & Pass (7 tests) & Passed \\
\hline
S4 & Admin RBAC & Non-admin rejected & Pass (18 tests) & Passed \\
\hline
S4 & Notifications lifecycle & CRUD + unread count & Pass (8 tests) & Passed \\
\hline
S1/S5 & Avatar CORP + health probe & Headers + DB/Redis ok & Pass (4 tests) & Passed \\
\hline
\end{tabularx}
\end{table}

\textbf{Analysis.} Each sprint increment in Table~\ref{tab:sprints} has at least
one passing integration-test group, so the chain from requirements through design
to executable verification stays intact. Administration (S4, 18 tests) and
document access (S2, 14 tests) form the largest groups, which matches how heavy
RBAC and visibility rules are in practice.
Figure~\ref{fig:ui-val-terminal-tests} shows the Vitest summary from the 7 July
2026 re-run (74/74, duration $\approx$41\,s).

\uiscreenshot{ui-val-terminal-tests.png}{Vitest integration suite output showing 74 tests passed across 10 files}{Integration test run}{fig:ui-val-terminal-tests}

\subsection{RAG Evaluation Results}

Quick evaluation without a configured Gemini API key classified all optimiser
topics as \texttt{general}, giving \textbf{15/34 PASS} on topic-label
cases. That drop is expected when the LLM optimiser cannot run. With API
credentials, the full harness writes JSON reports under
\texttt{apps/api/eval-reports/} for faithfulness and citation metrics.

\textbf{Interpretation.} Retrieval and serving paths can be checked without the
external API; semantic query optimisation and judged answer quality cannot.
The 15/34 figure is best treated as a \textbf{minimum regression gate}, not as a
score of production answer quality. Figure~\ref{fig:ui-val-rag-eval} records
the quick-mode harness output on 7 July 2026.

\uiscreenshot{ui-val-rag-eval.png}{RAG quick evaluation output showing 15 of 34 topic cases passed without Gemini API key}{RAG quick eval run}{fig:ui-val-rag-eval}

\subsection{Manual UX Validation}
\label{sec:manual-ux-validation}

Table~\ref{tab:manual-ux-tests} summarises the five manual browser scenarios.
The figures below supply visual evidence; matching implementation notes sit in
Chapter~\ref{ch:implementation} with complementary captions.

\begin{table}[H]
\centering
\caption{Manual UX validation summary}
\label{tab:manual-ux-tests}
\footnotesize
\renewcommand{\arraystretch}{1.35}
\setlength{\tabcolsep}{3pt}
\begin{tabularx}{\textwidth}{|>{\Centering\arraybackslash}p{0.7cm}|Y|>{\RaggedRight\arraybackslash}p{2.2cm}|>{\RaggedRight\arraybackslash}p{2.2cm}|>{\Centering\arraybackslash}p{1.2cm}|}
\hline
\textbf{Test} & \textbf{Objective} & \textbf{Expected Outcome} & \textbf{Observed Outcome} & \textbf{Status} \\
\hline
T1 & Reject invalid credentials without revealing account existence & Generic error message, no enumeration & Fig.~\ref{fig:ui-val-login-error} & Passed \\
\hline
T2 & Explain when document library access is disabled & Dedicated restricted page & Fig.~\ref{fig:ui-val-restricted} & Passed \\
\hline
T3 & Confirm ingest completes after upload & Processing status shown as complete on detail panel & Fig.~\ref{fig:ui-val-doc-ready} & Passed \\
\hline
T4 & Verify hybrid Ask returns cited answers & Answer with source citations & Fig.~\ref{fig:ui-val-ask-rag} & Passed \\
\hline
T5 & Show manager department overview & Dept-scoped team and documents & Fig.~\ref{fig:ui-val-mgr-dept} & Passed \\
\hline
\end{tabularx}
\end{table}

\paragraph{Authentication failure (T1).}
Invalid credentials produce a clear client-side error without revealing whether the
email exists, consistent with the auth integration suite
(Table~\ref{tab:test-summary}).

\uiscreenshot{ui-login-error.png}{Invalid email or password message after failed sign-in attempt}{Login error state}{fig:ui-val-login-error}

\paragraph{Feature restriction (T2).}
When document library access is disabled for a user, the API returns a feature
restriction response and the client redirects to a dedicated explanation page
(Figure~\ref{fig:seq-07}, Chapter~\ref{ch:architecture}).

\uiscreenshot{ui-access-restricted.png}{Dedicated page explaining that document library access is disabled}{Access restricted page}{fig:ui-val-restricted}

\paragraph{Ingest completion (T3).}
After upload and asynchronous processing, the file detail panel shows processing
status \textbf{READY}, confirming the ingest worker completed embedding and chunk
persistence.

\uiscreenshot{ui-document-details.png}{Document detail panel showing READY ingest status}{Ingest READY status}{fig:ui-val-doc-ready}

\paragraph{Hybrid Ask with citations (T4).}
The Ask interface returns a structured answer with confidence indicator and
clickable source citations grounded in retrieved chunks.

\uiscreenshot{ui-ask-rag-chat.png}{Ask response with structured answer and source citations}{Ask with citations}{fig:ui-val-ask-rag}

\paragraph{Manager department overview (T5).}
Managers view a read-only team directory and department-scoped document list for
manageable departments.

\uiscreenshot{ui-manager-department.png}{Manager department overview with team directory and documents}{Manager department view}{fig:ui-val-mgr-dept}

\section{Critical Discussion}
\label{sec:critical-discussion}

\subsection{Strengths}

End-to-end coherence is the clearest strength of the platform. Requirements remain
traceable, the architecture stays modular, the automated suite covers 74/74 cases,
and the documented RAG pipeline goes well past a thin chat demo: hybrid retrieval,
reranking, SSE, feedback, and an evaluation harness. Manual UX checks also show that
security and restriction flows are intelligible to users, not merely correct at the
API layer.

\subsection{Deviations from Targets}

A few targets were only partly met. The RAG judge KPI still lacks a full
faithfulness and relevance run, so it remains \textbf{Partial} in
Table~\ref{tab:kpi-results}. No formal load test was run for concurrent Ask
streams, leaving throughput unvalidated. The topic evaluation at 15/34 clears the
degraded threshold, yet it does not show production-grade query optimisation when
API credentials are absent.

\subsection{Limitations}

Several limits follow from the design choices. Latency, cost, and availability
all hinge on Google Gemini as the external LLM. Default storage is local disk
rather than cloud-native object storage. The automated judge is itself an LLM, so
human evaluation still has a place. Answer quality also climbs or falls with how
complete the ingest corpus is and how well chunking performs.

\subsection{Threats to Validity and Confidence Level}

Integration tests use a seeded database smaller than a production corpus. RAG eval
cases are code-defined and may miss some host-company domains. Manual UX tests
were run on a developer workstation.

\textbf{Confidence assessment.} Confidence is high for API contracts, RBAC,
document access, and notification lifecycles, backed by the 74/74 integration
suite. It is medium for RAG answer quality and citation faithfulness while the
eval harness remains partial without a full judge run. Confidence stays low for
concurrent load, production-scale corpus behaviour, and long-term operational
monitoring.

\subsection{Perspectives}

Natural next steps include SSO or LDAP integration, S3-compatible storage, on-prem
embedding models, and horizontal worker scaling. Multilingual UI, CI-integrated RAG
regression once API keys are available, and field deployment with user acceptance
testing on real organisational corpora would round out the present work.

\section{Chapter Conclusion}

This chapter checked the Knowledge Platform against the KPIs and sprint domains
from Chapter~\ref{ch:requirements}. Automated integration tests passed without
failure (74/74), ingest reliability cleared the 95\% target on the seeded corpus,
and five manual UX scenarios passed with screenshot evidence.
The RAG judge KPI is still partial until a full evaluation run with API credentials
is completed, and scale testing was not performed. The critical discussion places
high confidence in API-level correctness, medium confidence in RAG quality metrics,
and notes limits tied to the external LLM and the size of the test corpus. The
general conclusion weighs these outcomes against the objectives stated in the
introduction.
